% =========================================================
% PARTE 2: SIMULACIÓN Y GENERACIÓN DEL DATASET
% =========================================================
\section{Parte 2: Simulación y Generación del Dataset}

Para simular de manera estocástica y reproducible las 3 semanas de fin de cuatrimestre (21 días de operación continua), se dividió la jornada laboral (06:00 a 20:00 horas) en intervalos discretos de 15 minutos.

\subsection{Proceso Estocástico de Llegada de Alumnos (Poisson)}

La llegada de alumnos por intervalo de 15 minutos ($X_t$) en cada ruta se modela como una variable aleatoria Poisson:

\begin{equation}
	X_t \sim \text{Poisson}(\lambda_{\text{efectivo}})
\end{equation}

donde la tasa de llegada efectiva ($\lambda_{\text{efectivo}}$) se expresa según la tasa base ($\lambda_{\text{base}}$), la duración del intervalo ($\Delta t = 15$ min) y un multiplicador de hora pico ($M_{\text{pico}}$):

\begin{equation}
	\lambda_{\text{efectivo}} = \lambda_{\text{base}} \times \Delta t \times M_{\text{pico}}
\end{equation}

\begin{statbox}{Parámetros de Entrada por Ruta}
	\begin{itemize}
		\item \textbf{Puente de Tuzos:} $\lambda_{\text{base}} = \frac{1}{7}$ alumnos/min ($\approx 2.14$ pax/15 min). $M_{\text{pico}} = 1.6$. Tiempo medio de viaje: 20 min.
		\item \textbf{Central de Autobuses / Las Vías:} $\lambda_{\text{base}} = \frac{1}{3}$ alumnos/min ($\approx 5.0$ pax/15 min). $M_{\text{pico}} = 3.5$ ($< 1$ min por alumno). Tiempo medio de viaje: 27.5 min.
	\end{itemize}
\end{statbox}

\subsection{Estructura del Dataset Resultante}

El dataset simulado se exporta como archivo CSV (\texttt{data/demand\_dataset.csv}) con la siguiente especificación de columnas:

\begin{table}[H]
	\centering
	\caption{Estructura del Dataset Simulado Basado en Eventos de Despacho}
	\label{tab:dataset_schema}
	\begin{tabularx}{\textwidth}{>{\bfseries\color{purpleUPP}}l p{4.5cm} X}
		\toprule
		Categoría & Columna & Descripción Técnica \\
		\midrule
		Identificación & \texttt{id\_viaje} & Folio único correlativo de la corrida ($1, 2, 3 \dots$). \\
		& \texttt{fecha}, \texttt{ruta} & Fecha de operación y ruta (\textit{Puente Tuzos} vs \textit{Las Vías}). \\
		\addlinespace[0.4em]
		Despacho & \texttt{hora\_inicio\_espera} & Hora en que arribó el primer alumno a la unidad vacía. \\
		& \texttt{hora\_salida} & Estampa de tiempo efectiva de partida de la unidad. \\
		& \texttt{pasajeros} & Ocupación a bordo ($N = 18$ pax en política tradicional). \\
		& \texttt{es\_hora\_pico} & Indicador binario ($1 = \text{Hora Pico}$, $0 = \text{Hora Valle}$). \\
		\addlinespace[0.4em]
		Tiempos & \texttt{tiempo\_espera\_max\_min} & Minutos transcurridos para el primer pasajero (acumulado máximo). \\
		& \texttt{tiempo\_espera\_prom\_min} & Promedio de minutos esperados por los pasajeros del viaje. \\
		& \texttt{tiempo\_recorrido\_min} & Duración del trayecto a la UPP en minutos. \\
		& \texttt{hora\_llegada\_upp} & Estampa de tiempo de arribo a la universidad. \\
		\bottomrule
	\end{tabularx}
\end{table}

\newpage
